\chapter{How categories travel}
\label{ch:how-categories-travel}
\epigraph{It is a result of the theory I am urging that the use of metaphor in the expression of our theoretical beliefs is not a sign of the infancy of the discipline in question. It is rather a natural and appropriate response to certain features of the world.}{~-- Richard Boyd, \enquote{Metaphor and Theory Change} (1979)}

If categories are maintained clusters, and two maintained clusters share enough structure, then treating one as the other generates predictions. Some of those predictions hold. Some don't. The boundary between productive and unproductive extension tracks the boundary between shared and unshared mechanisms.

That claim has been implicit since Chapter~\ref{ch:kinds-without-essences}. This chapter makes it explicit, and in doing so addresses the tradition closest to HPC\ixsq{homeostatic property cluster!vs.\ cognitive linguistics} in spirit: cognitive linguistics\ixsq{cognitive linguistics}.

\section{The nearest tradition}
\label{sec:17:nearest-tradition}

Cognitive linguistics got several things right, and it got them right early.

\ixnq{Rosch, Eleanor H.}Rosch showed that categories have internal structure: typicality gradients, prototype effects, membership that comes in degrees of centrality \citep{rosch1975}. Chapter~\ref{ch:words-wont-hold-still} took that evidence seriously. \ixnq{Lakoff, George}Lakoff extended the insight to grammar, arguing that radial categories~-- the senses of \mention{over}\ixlq{over}, Dixon's Dyirbal\ixgq{Dyirbal} classifiers, the existential construction~-- show genuine clustering rather than definitional tidiness \citep{lakoff1987}. \ixnq{Lakoff, George}\ixnq{Johnson, Mark}\textcite{lakoff1980} demonstrated that cross-domain mappings are systematic, not ornamental: we don't just \emph{talk} about time in spatial terms; we \emph{reason} about it that way, and the mappings are productive enough to generate novel inferences. \ixnq{Bybee, Joan}\ixnq{Goldberg, Adele E.}\ixnq{Croft, William}Usage-based work grounded categories in actual use rather than innate specification \citep{bybee2010,goldberg1995constructions,croft2001}. And \ixnq{Langacker, Ronald W.}Langacker's Cognitive Grammar proposed a non-modular architecture in which usage events are the primitives and abstractions are emergent \citep{langacker1987foundations}.

These are genuine contributions. The programme documented clustering that essentialism couldn't see and took seriously the mechanisms~-- frequency, analogy, embodiment~-- that produce it. On both counts it anticipated the maintenance view this book defends: categories as mechanism-maintained profiles projectible relative to purposes.

But there's a gap, and it isn't empirical. The mechanisms are there~-- and in the usage-based wing, the ontological instinct is close to right. What's missing is the apparatus for saying when a maintained cluster actually \emph{projects}~-- when knowing the category lets you predict novel instances.

Prototype theory\ixsq{prototype} made real discoveries. \ixnq{Rosch, Eleanor}\ixnq{Mervis, Carolyn B.}\textcite{roschmervis1975} showed that prototypical members have the highest cue validity~-- their features best predict category membership~-- and that family resemblance structure (overlapping features, no single feature shared by all; \ixnq{Wittgenstein, Ludwig}\citealp{wittgenstein1953}) can be quantified. Basic-level effects~-- the finding that there's a privileged grain of categorization (\textsc{dog}, not \textsc{animal} or \textsc{poodle}) where informativeness peaks \citep{rosch1976}~-- identified where category structure is richest. These are substantial achievements, and they come with real explanatory content: cue validity tells you \emph{why} some members are more central than others. But it's distributional. It tells you which features co-occur; it doesn't tell you why they \emph{keep} co-occurring across speakers and centuries. Basic-level effects identify where the privileged level sits; they don't explain what maintains it there. The gap is in the maintenance-and-projection story: what keeps the clustering going, and what inductive work does the clustering underwrite?

If gradience were all there were, categories should drift: each generation learns from slightly different data, and small perturbations should accumulate. Chapter~\ref{ch:words-wont-hold-still} posed this as the stability problem. Embodiment\ixsq{embodiment} is cognitive linguistics' strongest reply: bodies persist across generations, so embodied concepts stay anchored. For spatial schemas and basic-level categories, the reply has force. But the grammatical categories this book examines~-- gender, deitality, count syntax, form--value couplings like \mention{done}/\mention{did}~-- aren't perceptually grounded. Their stability requires a different account. Cognitive linguistics proposed mechanisms that produce the gradience~-- entrenchment, frequency, analogical extension~-- but didn't develop a framework for predicting when those mechanisms keep a category stable, or when the resulting cluster projects~-- when knowing the category lets you predict things about novel instances that you couldn't predict without it.

Conceptual metaphor theory describes cross-domain mapping but doesn't explain what makes mapping \emph{possible}~-- what property of categories enables them to travel. \textcite{lakoff1980} showed that we systematically map spatial structure onto temporal reasoning. But not every domain maps onto every other. Why do some mappings work and others don't? What distinguishes a productive mapping from a dead one? Image schemas\ixsq{image schema} and experiential gestalts were serious attempts at an answer: experientially basic structures that serve as source domains. They identify what gets mapped. But they don't predict which new mappings should be productive, or explain why structural similarity licences inferential transfer in some cases and not others.

The HPC framework doesn't replace cognitive linguistics. It grounds it. Prototype effects are real because mechanisms produce and sustain them (Chapters~\ref{ch:stabilizers}--\ref{ch:failure-modes}). Radial categories are genuine clusters because the same maintenance processes~-- entrenchment, alignment, functional pressure~-- keep the extensions anchored to the core. And cross-domain mapping works when it works because of something specific about the categories involved: they share mechanism-sustained property overlap. Where the mechanisms diverge, the mapping breaks down. That's a testable claim, not a restatement.

Prototype theory identified the pattern; HPC explains why the pattern holds.

What the cognitive linguists lacked wasn't data, ambition, or even the right ontological instinct. Bybee's entrenchment, Goldberg's constructional schemas, Tomasello's usage events all point toward a constitutive role for mechanisms~-- categories \emph{are} their usage patterns, not merely described by them. What was missing was the diagnostic toolkit: projectibility criteria that say when a cluster supports induction (Chapter~\ref{ch:projectibility}), failure modes that diagnose when it doesn't (Chapter~\ref{ch:failure-modes}), and field-relativity that explains why the same phenomena yield different categories for different purposes.

The sharpest gap is counting. Prototype theory can tell you a category has internal structure~-- radial extensions, sub-categories, typicality gradients. What it can't tell you is whether that structure is one category with sub-parts or two categories whose extensions overlap. Chapter~\ref{ch:definiteness-and-deitality} is exactly this ambiguity resolved: what looks like one fuzzy \mention{definiteness} with peripheral members turns out to be two clusters~-- semantic identifiability and formal deitality~-- maintained by partially independent mechanisms, decoupling where the mechanisms diverge. The proper noun/proper name distinction (§\ref{sec:7:proper-nouns-names}) makes the same move in miniature: same forms, different maintenance, different projections. HPC resolves the ambiguity because it asks what maintains each part of the cluster. Prototype theory can't, because it has no maintenance story to partition with. You can hold the right ontology without having the means to test it.

This chapter closes that gap, and the closing has a payoff: cross-domain extension falls out of the framework for free~-- as a consequence of what categories are.

\section{Theory-constitutive metaphor and cross-domain extension}
\label{sec:17:tcm}

Richard Boyd made an argument in 1979 that the philosophy of science absorbed and linguistics largely ignored.

Metaphors in science, Boyd argued, aren't pedagogical aids~-- scaffolding that can be replaced by literal paraphrase once the theory matures. Some metaphors are \term{theory-constitutive\ixsq{theory-constitutive metaphor}}: they express causal relationships between categories that no literal description yet captures \citep{boyd1979}. When a biologist calls natural selection a \enquote{pressure}, the metaphor isn't ornamental. It constitutes a claim about the dynamics of fitness landscapes that no available literal vocabulary can state as precisely. The metaphor succeeds not because source and target share surface features but because they share \emph{mechanism-sustained property overlap}. Pressure in physics and selection in biology are both maintained processes that produce directional change; the shared structure licences inferential transfer.

This is more principled than Lakoff and Johnson's account of mapping. \textcite{lakoff1980} demonstrated that cross-domain mappings are systematic~-- we don't map randomly from source to target. Their framework tracks which mappings are conventional and which are novel, and it distinguishes dead metaphors from live ones. But it doesn't predict which \emph{new} mappings should be productive, or explain why some domains map onto each other and others don't. Image schemas provide a partial answer (spatial structure is experientially basic and therefore available as a source), but they don't explain why TIME IS MONEY works while TIME IS FURNITURE doesn't, even though both furniture and money are culturally familiar domains with rich internal structure.

Boyd's answer: the mapping works to the extent that source and target are HPC kinds whose partially overlapping mechanisms sustain projectible inference. We budget, spend, save, and waste both time and money~-- Lakoff and Johnson's inventory is well known~-- but the HPC account explains what sustains those correspondences: both time and money are maintained resources, and the resource-management mechanisms overlap. Invest time and you can predict returns; waste it and you bear opportunity cost. The inferences project because the mechanisms are shared. \textsc{time is furniture} fails the same test. You can arrange both furniture and an afternoon, but nothing else transfers: well-arranged furniture stays put; a well-arranged afternoon doesn't. The overlap is in surface features, not in the causal architecture that maintains them or the inferences that architecture underwrites.

Define \term{cross-domain extension}\ixsq{cross-domain extension} as follows: when two HPC kinds share partial mechanism-sustained property overlap, inferences licensed by one can be extended to the other. The extension is productive to the extent that the shared properties are maintained by shared or analogous mechanisms. This is what metaphor \emph{is}: a consequence of what categories are. If categories are maintained clusters, and two maintained clusters share enough structure, then treating one as the other generates predictions. Some hold (the metaphor works); some don't (the metaphor breaks down). The boundary between productive and unproductive extension tracks the boundary between shared and unshared mechanisms.

Where maintenance overlaps, projection overlaps.

This yields a prediction: wherever you have HPC kinds, you should find cross-domain extension. The prediction is testable. Cross-domain extension should be more productive between kinds that share more mechanism-sustained structure, and it should degrade when the shared mechanisms are disrupted.

\ixnq{Noyes, Alexander}\ixnq{Keil, Frank C.}\textcite{noyes2019} provide experimental teeth for this philosophical claim. Generic statements (\mention{Dogs follow trails}\ixlq{dog}) transmit what Noyes and Keil call the \term{kind assumption}~-- the belief that a category has rich internal structure supporting inductive generalization~-- without necessarily transmitting an \term{essence assumption} about hidden determining natures. In their experiments, generic framing produced a large effect on kind-related beliefs but only a small one on essence-related beliefs ($d = 1.09$ vs.\ $d = 0.37$); when the attributed properties were cultural rather than biological, the essence effect vanished while the kind effect held. The kind assumption also generalized across animals, artifacts, and social categories, confirming that the cognitive recognition of rich structure is domain-general even when the maintaining mechanisms differ. Categories can be inductively rich without being essential, and the apparatus for recognizing that richness doesn't respect domain boundaries.

\ixnq{Gentner, Dedre}\textcite{gentner1983} formalized a closely related insight with her structure-mapping theory: analogies are productive when they preserve relational structure, not when they match surface features. HPC adds the missing piece: what makes relational structure \emph{available} for mapping is that it's maintained by mechanisms. Dead metaphors die when the mechanisms diverge; live metaphors live because shared maintenance keeps the relational structure aligned.

Literary metaphor\ixsq{literary metaphor} inverts the logic. \ixnq{Dickinson, Emily}Dickinson's \enquote{Hope is the thing with feathers} maps a richly structured source onto a richly structured target, and no shared mechanisms sustain the mapping. That's the resource, not the limitation. The reader keeps finding fits~-- persistence, fragility, unbidden arrival~-- and misfits. Hope doesn't moult. If mechanisms sustained the mapping, it would conventionalize and stop being interesting. The poet's art is precise selection at the mechanism boundary: activating exactly the right overlap and refusing to extend it further. It's mechanism support, not frequency, that conventionalizes a metaphor \citep[see][for the full argument]{reynolds2026metaphor}.

\section{How categories travel: the trail}
\label{sec:17:trail}

To see cross-domain extension at work, consider a category that predates language by tens of millions of years.

Dogs navigate primarily by olfaction. When a person walks across a field, they deposit a scent trail: volatile organic compounds from skin, sweat glands, and the soles of shoes, supplemented by crushed vegetation and disturbed soil. The trail has structure. \ixnq{Thesen, A.}\textcite{thesen1993} showed that trained dogs can determine the direction of a human trail from as few as five footprints, using the concentration gradient~-- newer prints carry stronger scent. \ixnq{Hepper, Peter G.}\ixnq{Wells, Deborah L.}\textcite{hepper2005} replicated the finding and added that dogs discriminate trail direction even when the prints are only seconds apart, relying on decay asymmetries too subtle for human detection.

What makes a trail a trail? Not any single property. A trail has directionality, recency, a concentration profile, a substrate type, and a source identity~-- and these cluster together, not by logical necessity but because mechanisms keep them co-occurring. The dog's own scent marking reinforces the trail; route reuse deepens it; the substrate's physical persistence (moisture, temperature, porosity) modulates how long the cluster holds together. \textsc{Trail}\ixsq{trail} is an HPC category, maintained the same way grammatical categories are, just with different mechanisms. And the same removal test applies: dry the substrate, eliminate reuse, suppress scent marking, and the cluster degrades.

A tracking dog trained on grass can follow the same person's trail across concrete, carpet, and leaf litter. That looks unremarkable~-- it's just what tracking dogs do. But notice what's changing underneath. Scent compounds persist differently on moist grass, non-porous concrete, and absorbent carpet; concentration gradients and decay dynamics shift with every surface. Some properties of the \textsc{trail} cluster carry across (source identity, directionality); others don't (concentration profile, substrate interactions). The dog holds onto what transfers and adjusts for the rest.

Now a harder case. My own dog follows the painted lines on soccer fields. The painted line shares visual-geometric trail properties~-- linear extent, bounded width, high contrast against the surrounding surface~-- but none of the original mechanisms are operating. Nobody walked it into existence; no scent trail reinforces it; no foot traffic deepens it. But the dog recruits trail-following behaviour on the basis of the partial property overlap that remains. The extension is productive (the dog follows the line) but limited (there's nothing to track at the end). That limitation is informative: it marks exactly where the source cluster's properties stop transferring \citep[see][for the fuller argument and confound analysis]{reynolds2026metaphor}.

The two cases differ in degree, not in kind. In both, the dog extends \textsc{trail} beyond the conditions that originally maintained it. What changes is how much of the mechanism suite carries over. Call two instances part of the same \term{functional domain}\ixsq{functional domain} when the same mechanisms maintain their property clusters. Cross-domain extension occurs when the suite changes~-- partially for the tracking dog on concrete, almost completely for the dog on the painted line.

\ixnq{Fagan, William F.}\textcite{fagan2025} push the point further. Studying GPS-tracked movement data across species, they find that wild canids show significantly stronger route reuse than felids. Wolves, coyotes, and African wild dogs return to established trails; big cats don't. The difference tracks mechanism density\ixsq{mechanism!density}: canids are social scent-markers that reinforce trails through collective use, while felids are largely solitary hunters whose movement patterns leave shallower traces. The \textsc{trail} category is more deeply entrenched in canid cognition~-- more properties cluster, more mechanisms maintain the clustering, and the category supports richer extension.

The phylogenetic\ixsq{phylogenetic depth} point follows, but it needs precision. The claim isn't that dogs think metaphorically~-- that would be semiotic overreach. The claim is narrower: the capacity to extend a category when the maintaining mechanisms partially change is evolutionarily old. Dogs have it. What humans add is semiotic amplification~-- signs detachable from immediate context, interpretants\ixsq{interpretant} that chain across domains, habits transmittable through culture~-- that lets the same basic capacity operate between conceptual domains, not just physical substrates. A tracking dog adjusting to concrete and a speaker mapping \textsc{time is money} are not doing the same thing. But the gap between them is one of semiotic reach, not of kind. Language doesn't create cross-domain extension. It extends how far a category can travel.

Large language models\ixsq{large language models} sharpen the contrast. They reliably detect structural correspondences across domains~-- often enough to make their analogies feel apt~-- but their selectivity is largely extrinsic, imposed by prompting and post-training rather than anchored in endogenous goals and world-coupled costs. Absent a task-framed stopping condition, they will elaborate correspondences that a situated agent would treat as dead ends. They follow any line, including the painted ones. Endogenous goals are what let the dog be selective~-- not every overlap is worth pursuing, and a dog with a scent to track won't waste time on paint. Semiotic amplification is what lets the human extend further than the dog can, projecting categories across conceptual domains rather than just physical substrates.

\section{The Boyd circle}
\label{sec:17:boyd-circle}

The argument is now reflexive, and the reflexivity is a feature.

Start the circle:

\begin{enumerate}
\item \textcite{boyd1979} argued that theory-constitutive metaphors work because source and target are both HPC kinds sharing partial property overlap.
\item \textcite{boyd1991} developed the HPC framework for species in biology: organisms cluster because reproductive, developmental, and ecological mechanisms keep properties co-occurring.
\item This book applies HPC~-- originally from species~-- to linguistic categories, with Part~III centred on grammatical categories. That application is itself a cross-domain extension on partial property overlap: species and linguistic categories share mechanism-sustained clustering, projectibility, and boundary behaviour, even though the specific mechanisms differ.
\item The book now predicts that HPC kinds should support cross-domain extension wherever the mechanism-sustained overlap is rich enough.
\item Cross-domain extension is what metaphor \emph{is}.
\item Which is what Boyd was explaining in step 1.
\end{enumerate}

This isn't vicious circularity. The theory doesn't presuppose its own conclusion. It predicts a capacity~-- cross-domain extension~-- that the theory itself exemplifies. The reflexivity is a test. If the framework couldn't account for its own method of explanation, that would be a problem: a theory of categories that can't explain why it was able to travel from biology to linguistics would be undermining its own explanatory claims. It can. The self-application passes because the property overlap between biological species and linguistic categories is genuine and mechanism-sustained~-- a structural analogy grounded in shared causal architecture (maintained clustering, projectible properties, homeostatic boundaries).

The Peircean\ixsq{Peirce, Charles Sanders} vocabulary from Chapter~\ref{ch:projectibility} sharpens the point. The cable metaphor \citep[CP~5.265]{peirce1931cp} is itself a theory-constitutive metaphor in Boyd's sense. Peirce used the physical properties of cables~-- multiple partially independent strands whose collective strength exceeds any single strand~-- to constitute a claim about epistemic justification. The metaphor works because cables and arguments share HPC-like structure: maintained by distributed redundancy, resistant to local failure, degrading gracefully rather than catastrophically. This book's central move~-- applying HPC from biological species to linguistic categories~-- is a TCM in exactly the same sense. It works because species and linguistic categories share real causal architecture (maintained clustering, projectible properties, homeostatic boundaries), not because someone stipulated the analogy.

\ixnq{Hofstadter, Douglas R.}\ixnq{Sander, Emmanuel}\textcite{hofstadter2013} argue that analogy is the \enquote{fuel and fire of thinking}~-- that every act of categorization is an act of analogy, mapping new instances onto existing categories by structural similarity. The HPC framework specifies what that structural similarity consists of: shared mechanism-sustained properties. Hofstadter and Sander are right that analogy is pervasive. What they leave unspecified is what makes one analogy productive and another empty. A productive analogy is one that projects~-- where learning something about the source lets you make reliable predictions about the target. And projection, as the rest of this book argues, is sustained by shared causal architecture: overlapping mechanisms that keep properties clustering in both domains. The quality of an analogy tracks that shared architecture, not the vividness of the surface match.

The explanatory coherence runs deeper than self-application. The framework that explains what linguistic categories are also explains how theories about linguistic categories get built~-- through the same cross-domain extension it predicts. This is a structural consequence of the theory's content: if categories are maintained clusters that support inductive extension, and if theories are themselves maintained clusters of concepts and commitments, then the capacity for theory-building is a special case of the capacity for categorization. The circle closes because it was never really a circle. It's a single capacity~-- forming maintained clusters that project~-- operating at different scales.

\section{What this means}
\label{sec:17:what-this-means}

For cognitive linguistics, five consequences.

First, prototype effects are mechanistically grounded, not primitive. When Rosch found that robins are better birds than penguins, she identified a real gradient and proposed cue validity as its explanation~-- a genuine one. The HPC framework adds the maintenance story: the same processes~-- entrenchment, acquisition bias, functional pressure~-- keep regenerating the cue-validity structure in each speaker and across generations. Prototype theory quantified the gradient and explained its shape. Homeostasis explains why the shape persists and why it projects~-- why a new speaker's typicality judgments are predictable from the old ones.

Second, cross-domain mappings work because of shared mechanism-sustained structure. Image schemas\ixsq{image schema} may well be real psychological structures, and nothing here denies their existence. But they describe the \emph{shape} of the overlap without explaining its \emph{source}. The HPC account says what source and target have to share for the mapping to project: overlapping causal architecture. A mapping between two HPC kinds with rich overlapping mechanisms will generate novel inferences that hold. Where the shared architecture runs deep enough, the transfer goes both ways: properties of each domain become defeasible properties of the other, and learning about the target feeds back into understanding the source. The points where properties stop transferring are diagnostic of where the maintenance regimes diverge. A mapping where the overlap is accidental~-- shared features without shared maintenance~-- will be shallow, generating surface analogies that break under load. The methodological upshot: stop cataloguing the mapping and start asking whether it projects. Which mechanisms does the source domain share with the target? Where do they diverge? That's where the metaphor stops generating true predictions, and that's where the most informative data lives.

Third, the framework is self-applying: it predicts the very capacity it uses. This adds a falsification condition that cognitive linguistics lacks. If productive cross-domain extension~-- the kind that generates true predictions in the target~-- doesn't track shared causal architecture, the framework loses explanatory scope. Rich metaphors arising between domains with no overlapping mechanisms, or impoverished metaphors persisting between domains with substantial shared maintenance, would both count. The prediction is on the table.

Fourth, the failure modes from Chapter~\ref{ch:failure-modes} apply. If the source is a fat class\ixsq{fat class}~-- multiple kinds conflated~-- the metaphor will be ambiguous, working for one sub-cluster but not another. If it's thin\ixsq{thin class}, the metaphor is shallow: not enough structure to sustain transfer. A negative class\ixsq{negative class} can't source a metaphor at all~-- no positive architecture, no structure to map. The diagnostics that tell you whether a category is real tell you whether a metaphor built on it will work.

Fifth, the framework can count categories. Prototype theory can see internal structure~-- sub-categories, peripheral members, typicality gradients~-- but it can't tell you whether that structure is one category with sub-parts or two categories whose extensions overlap. HPC can, because different maintenance regimes generate different projections. Definiteness and deitality project differently; proper nouns and proper names project differently. Treating either pair as one category mispredicts at every boundary.

A methodological caution about the framework's own architecture. The HPC triad~-- clustering, maintenance, projection~-- is tested with three independent diagnostics, but in practice the three legs strongly covary. Maintained clusters almost always project; clusters that project almost always turn out to be maintained. That covariation isn't a coincidence the theory ignores; it's what the theory predicts. Mechanisms sustain clusters \emph{because} doing so supports reliable induction. But the covariation means that most of the time, asking \enquote{is this maintained?} and \enquote{does this project?} will give the same answer. The cases where they come apart~-- the defeat conditions in Appendix~\ref{app:predictions} and the failure modes in Chapter~\ref{ch:failure-modes}~-- are where the distinction between maintenance and projection does its sharpest work. If those cases never arise empirically, the triad is philosophical scaffolding rather than empirical architecture. That would be worth knowing.

For the book's programme, this chapter adds two things. The constructive argument (Parts II and III) showed that linguistic categories are maintained clusters, projectible relative to purposes. Chapter~\ref{ch:what-changes} demands predictions and defeat conditions. This section adds one more prediction: productive cross-domain extension~-- the kind that generates true predictions about the target~-- should track shared causal architecture. And it provides one more piece of evidence: the theory's own history~-- its journey from biology to linguistics to a theory of metaphor~-- instantiates the capacity it describes. If that's a coincidence, it's a remarkably well-structured one.

That history also explains why Chapter~\ref{ch:what-changes} can demand the predictions it demands. The application of HPC from biological species to linguistic categories isn't an analogy that motivated the book and can now be discarded. It's a theory-constitutive metaphor in Boyd's sense: the shared causal architecture between species and linguistic categories~-- maintained clustering, projectible properties, patterned boundary behaviour, degradation when mechanisms weaken~-- is what licences the inferential transfer. The predictions tested there (asymmetric fraying, sticky-lock drift, phase transition in genesis) are generated by that transfer: they hold for biological species, the property overlap is mechanism-sustained, and so the framework predicts they should hold for linguistic categories too. If the overlap were merely surface resemblance~-- if species and linguistic categories shared features without sharing maintenance architecture~-- the predictions would be decorative, not load-bearing. They carry weight because the metaphor is constitutive.
